\documentclass[a4paper,12pt]{article}
\usepackage[utf8]{inputenc}
\usepackage[ngerman]{babel}
\usepackage{amsmath,amssymb}
\usepackage{geometry}
\geometry{margin=2.5cm}

\title{Aufgabe 7(a) -- Gauß- und Gauß-Jordan-Elimination}
\author{}
\date{}

\begin{document}
\maketitle

\section*{Gleichungssystem}

\[
\begin{cases}
4x + 7y = -1 \\
y - 2z = -5 \\
2x + y = -3 \\
x + y + z = 2
\end{cases}
\]

\section*{Erweiterte Koeffizientenmatrix}

\[
\left(\begin{array}{ccc|c}
4 & 7 & 0 & -1 \\
0 & 1 & -2 & -5 \\
2 & 1 & 0 & -3 \\
1 & 1 & 1 & 2
\end{array}\right)
\]

\textbf{Tipp:} Wir tauschen Zeile 1 und Zeile 4, damit wir eine 1 links oben haben -- das macht die Rechnung einfacher.

\section*{Gauß-Elimination (Stufenform)}

\textbf{Schritt 1:} $Z_1 \leftrightarrow Z_4$

\[
\left(\begin{array}{ccc|c}
1 & 1 & 1 & 2 \\
0 & 1 & -2 & -5 \\
2 & 1 & 0 & -3 \\
4 & 7 & 0 & -1
\end{array}\right)
\]

\textbf{Schritt 2:} $Z_3 \leftarrow Z_3 - 2Z_1$ und $Z_4 \leftarrow Z_4 - 4Z_1$

\[
\left(\begin{array}{ccc|c}
1 & 1 & 1 & 2 \\
0 & 1 & -2 & -5 \\
0 & -1 & -2 & -7 \\
0 & 3 & -4 & -9
\end{array}\right)
\]

\textbf{Schritt 3:} $Z_3 \leftarrow Z_3 + Z_2$ und $Z_4 \leftarrow Z_4 - 3Z_2$

\[
\left(\begin{array}{ccc|c}
1 & 1 & 1 & 2 \\
0 & 1 & -2 & -5 \\
0 & 0 & -4 & -12 \\
0 & 0 & 2 & 6
\end{array}\right)
\]

\textbf{Schritt 4:} $Z_4 \leftarrow 2Z_4 + Z_3$

\[
\left(\begin{array}{ccc|c}
1 & 1 & 1 & 2 \\
0 & 1 & -2 & -5 \\
0 & 0 & -4 & -12 \\
0 & 0 & 0 & 0
\end{array}\right)
\]

\textbf{Ergebnis:} Die Stufenform hat 3 Pivot-Elemente und keine Widerspruchszeile $\Rightarrow$ das System ist \textbf{konsistent} mit einer \textbf{eindeutigen Lösung}.

\subsection*{Rückwärtseinsetzen}

Aus Zeile 3: $-4z = -12 \Rightarrow z = 3$

Aus Zeile 2: $y - 2z = -5 \Rightarrow y = -5 + 2 \cdot 3 = 1$

Aus Zeile 1: $x + y + z = 2 \Rightarrow x = 2 - 1 - 3 = -2$

\[
\boxed{(x, y, z) = (-2, 1, 3)}
\]

\section*{Gauß-Jordan-Elimination (reduzierte Stufenform)}

Wir starten von der Stufenform und eliminieren nach oben.

\textbf{Schritt 5:} $Z_3 \leftarrow -\frac{1}{4} Z_3$ (Pivot normieren)

\[
\left(\begin{array}{ccc|c}
1 & 1 & 1 & 2 \\
0 & 1 & -2 & -5 \\
0 & 0 & 1 & 3 \\
0 & 0 & 0 & 0
\end{array}\right)
\]

\textbf{Schritt 6:} $Z_2 \leftarrow Z_2 + 2Z_3$ und $Z_1 \leftarrow Z_1 - Z_3$

\[
\left(\begin{array}{ccc|c}
1 & 1 & 0 & -1 \\
0 & 1 & 0 & 1 \\
0 & 0 & 1 & 3 \\
0 & 0 & 0 & 0
\end{array}\right)
\]

\textbf{Schritt 7:} $Z_1 \leftarrow Z_1 - Z_2$

\[
\left(\begin{array}{ccc|c}
1 & 0 & 0 & -2 \\
0 & 1 & 0 & 1 \\
0 & 0 & 1 & 3 \\
0 & 0 & 0 & 0
\end{array}\right)
\]

Dies ist die \textbf{reduzierte Zeilenstufenform} (RREF). Man liest direkt ab:

\[
\boxed{(x, y, z) = (-2, 1, 3)}
\]

\section*{Probe}

\begin{align*}
4(-2) + 7(1) &= -8 + 7 = -1 \checkmark \\
1 - 2(3) &= 1 - 6 = -5 \checkmark \\
2(-2) + 1 &= -4 + 1 = -3 \checkmark \\
-2 + 1 + 3 &= 2 \checkmark
\end{align*}

\end{document}
